\chapter{Validazione Sperimentale e Discussione}
\label{chap:validazione}

In questo capitolo vengono presentati i risultati ottenuti testando la 
pipeline decisionale proposta nel Capitolo~\ref{chap:metodologia}. L'obiettivo 
principale è verificare se la modellazione geometrica con ellissi e la pianificazione 
dinamica migliorino concretamente la scelta della sequenza di raccolta 
rispetto ai metodi tradizionali basati su \textit{bounding box}.

Dapprima vengono descritte le immagini di test acquisite in serra e i 
parametri operativi della pipeline. Successivamente, l'efficacia del metodo 
viene valutata confrontando le decisioni dell'algoritmo con la baseline a 
riquadri rigidi e con le scelte di riferimento annotate da operatori umani. 
Infine, l'analisi di alcuni casi studio reali consente di esaminare da vicino 
come il sistema gestisce lo sblocco dei frutti occlusi e quali sono i 
principali limiti riscontrati.

\section{Configurazione Sperimentale e Dataset di Test}
\label{sec:setup_sperimentale}

La correttezza della pianificazione dipende direttamente dalla qualità 
delle maschere fornite dal modello di visione. Contorni rumorosi o 
frammentati alterano il calcolo dell'involucro convesso e dell'ellisse 
associata a ciascun frutto, inducendo stime errate delle occlusioni 
tra pomodori adiacenti. Definire con rigore le condizioni di acquisizione, 
la risoluzione di inferenza e le soglie geometriche costituisce quindi 
il presupposto necessario per isolare l'efficacia del planner e 
garantirne la riproducibilità.


\subsection{Dataset e Immagini di Test}
\label{subsec:dataset}

La sperimentazione è stata condotta a partire da un dataset acquisito in 
ambiente di serra, composto complessivamente da $1803$ immagini con 
risoluzione nativa di $432 \times 432$ pixel. Le annotazioni poligonali 
comprendono quattro classi: frutti (\textit{tomato}), fusto principale 
(\textit{MainStem}), peduncoli (\textit{peduncle}) e piccioli 
(\textit{petiole}). Ai fini della raccolta, la pipeline sviluppata elabora 
esclusivamente le maschere appartenenti alla classe \textit{tomato}.

Per la validazione della pipeline decisionale e del task planning, è stato 
selezionato un sottoinsieme mirato di $49$ immagini. Questi fotogrammi sono 
stati scelti specificamente per la presenza di grappoli densi con marcate 
occlusioni reciproche, sovrapposizioni parziali e ombreggiature fogliari, 
costituendo il banco di prova ideale per verificare e confrontare l'efficacia 
delle diverse strategie di prelievo.


\subsection{Configurazione del Modulo Percettivo (YOLOv11-seg)}
\label{subsec:modulo_percettivo}

La segmentazione delle istanze costituisce il primo stadio della pipeline decisionale. 
In questo lavoro si è scelto di utilizzare il modello \textbf{YOLO11-x} (Extra-Large) 
con risoluzione di inferenza a $\text{imgsz} = 800$ pixel, configurazione già ottimizzata 
e convalidata nelle precedenti attività di laboratorio~\cite{collega_yolo} con un 
$\text{mAP@50-95}$ del $40.8\%$ e una precisione del $79.0\%$.

Il presente lavoro si concentra infatti sulla fase successiva al rilevamento: partendo 
dalle maschere binarie fornite dalla rete, l'obiettivo è determinare una sequenza di 
raccolta ordinata, efficiente e priva di collisioni. Per questo motivo la soglia di 
confidenza è stata mantenuta a $\text{conf} = 0.30$, così da massimizzare la sensibilità 
del modello anche sui pomodori parzialmente nascosti e lasciare alle successive fasi 
geometriche il compito di filtrare le maschere non affidabili.

\subsection{Ambiente di Esecuzione e Parametri Operativi}
\label{subsec:parametri_operativi}

Tutti i moduli dell'algoritmo sono stati implementati in linguaggio Python~3.12, 
sfruttando le librerie PyTorch per l'inferenza del modello di deep learning e 
OpenCV per le primitive algebriche e geometriche (Convex Hull, fitting ellittico, 
trasformazioni nello spazio colore HSV e operazioni morfologiche binarie).

A fini di tracciabilità e riproducibilità sperimentale, in 
Tabella~\ref{tab:iperparametri_pipeline} vengono riepilogati tutti i parametri 
operativi fissati per i test della pipeline decisionale, coerentemente con le 
formulazioni matematiche dedotte nel Capitolo~\ref{chap:metodologia}.

\begin{table}[htbp]
    \centering
    \small
    \setlength{\tabcolsep}{5pt}
    \caption{Riepilogo dei parametri e delle soglie operative della pipeline decisionale.}
    \label{tab:iperparametri_pipeline}
    \begin{tabular}{lll p{6.0cm}}
        \toprule
        \textbf{Parametro} & \textbf{Simbolo} & \textbf{Valore} & \textbf{Ruolo Operativo} \\
        \midrule
        Area minima ammissibile & $A_{\min}$ & $1500$ px & Esclusione preliminare di falsi positivi e sfondo \\
        Soglia di deduplicazione & $\tau_{\text{dedup}}$ & $0.70$ & Mask-IoU per l'eliminazione di predizioni duplicate \\
        Dimensione del pool & $K$ & $8$ & Numero massimo di candidati ammessi alla Fase 3 \\
        Soglia di maturità & $\tau_{\text{mat}}$ & $0.50$ & Indice HSV $M_i$: discrimina maturi da ostacoli acerbi \\
        Soglia di occlusione & $\tau_{\text{occ}}$ & $0.15$ & Rapporto di overlap tra ellissi $\rho_{ij}$ per contatto \\
        Penalità di occlusione & $\gamma_{\text{occ}}$ & $0.70$ & Fattore applicato a frutti frontalmente occlusi \\
        Bonus di sblocco & $\beta_{\text{unlock}}$ & $1.20$ & Moltiplicatore assegnato a frutti liberati da un prelievo \\
        Lunghezza sequenza & $N_{\max}$ & $3$ & Numero di target pianificati per ciclo di picking \\
        \bottomrule
    \end{tabular}
\end{table}

\section{Analisi Comparativa delle Strategie di Picking}
\label{sec:analisi_comparativa}

La validazione sperimentale richiede un confronto quantitativo diretto con le 
metodologie convenzionali di selezione dei target. L'obiettivo è verificare 
se la modellazione mediante ellissi orientate e la ripianificazione dinamica 
offrano un incremento misurabile nell'accuratezza della sequenza di raccolta 
rispetto all'impiego dei tradizionali riquadri rigidi (\textit{bounding box}).

\subsection{Metodologia di Benchmark e Criteri di Valutazione}
\label{subsec:metodologia_benchmark}

Per quantificare l'impatto della modellazione proposta, è stato condotto un 
confronto sperimentale a tre vie su immagini reali acquisite in serra:
\begin{enumerate}
    \item \textbf{Ground Truth Umano}: sequenza di raccolta ideale (fino a 
    tre frutti: $H_1, H_2, H_3$) annotata manualmente da operatori su un 
    sottoinsieme di $21$ scene a elevata densità di grappolo, selezionata 
    valutando l'accessibilità visiva e l'assenza di ostacoli frontali.
    \item \textbf{Baseline Bounding Box}: implementazione di 
    riferimento~\cite{collega_yolo} che modella i pomodori tramite riquadri 
    rigidi allineati agli assi (AABB) o cerchi isotropi, determinando le 
    occlusioni e la sequenza in modo statico.
    \item \textbf{Planner Ellittico Proposto}: la pipeline descritta nel 
    Capitolo~\ref{chap:metodologia}, basata su fitting di ellissi orientate 
    ($a_i, b_i, \theta_i$), calcolo dell'Area Ownership per l'ordinamento in 
    profondità e ripianificazione dinamica con bonus di sblocco 
    $\beta_{\text{unlock}}$.
\end{enumerate}

La valutazione si concentra su due criteri principali: la concordanza sulla 
scelta del primo frutto da raccogliere (Rank~1), che rappresenta la decisione 
più critica per prevenire collisioni nell'approccio iniziale del robot, e 
l'aderenza della sequenza completa ($N_{\max} = 3$) rispetto alla strategia 
umana.


\subsection{Risultati Quantitativi Globali}
\label{subsec:risultati_globali}

L'analisi condotta sull'intero set di test ($43$ scene valide) ha evidenziato 
la seguente distribuzione dei comportamenti decisionali:
\begin{itemize}
    \item In $27$ scene su $43$ ($62.8\%$), i due algoritmi producono la 
    medesima sequenza di prelievo. Questo risultato attesta che in grappoli 
    sparsi o con sovrapposizioni minime la modellazione ellittica preserva la 
    coerenza con la baseline, senza introdurre distorsioni.
    \item In $12$ scene ($27.9\%$), l'algoritmo ellittico diverge dalla 
    baseline, modificando il target selezionato come Rank~1 oppure invertendo 
    l'ordine di prelievo dei frutti successivi.
    \item In $4$ scene ($9.3\%$), entrambi i metodi segnalano l'assenza di 
    frutti maturi accessibili, comportandosi correttamente come meccanismo di 
    sicurezza (\textit{safe-fail}) in presenza di soli pomodori verdi o 
    totalmente occlusi.
\end{itemize}

Per valutare la bontà delle decisioni prese dall'algoritmo, si è confrontata 
la sequenza calcolata con quella annotata manualmente dagli operatori sulle 
$21$ scene ad alta densità. La Tabella~\ref{tab:benchmark_ground_truth} 
riassume la percentuale di concordanza con il riferimento umano per ciascuna 
delle tre posizioni della sequenza di raccolta ($N_{\max} = 3$).

\begin{table}[htbp]
    \centering
    \small
    \caption{Tasso di concordanza con il Ground Truth umano lungo la sequenza 
    di raccolta.}
    \label{tab:benchmark_ground_truth}
    \begin{tabular}{lcccc}
        \toprule
        \textbf{Posizione Target} & \textbf{Scene Valutate} & 
        \textbf{Baseline BBox} & \textbf{Planner Ellittico} & 
        \textbf{Variazione ($\Delta$)} \\
        \midrule
        Rank 1 & 21 & 17 / 21 (81.0\%) & 17 / 21 (81.0\%) & $+0.0\%$ \\
        Rank 2 & 13 & 5 / 13 (38.5\%)  & 6 / 13 (46.2\%)  & $+7.7\%$ \\
        Rank 3 & 10 & 2 / 10 (20.0\%)  & 4 / 10 (40.0\%)  & $+20.0\%$ \\
        \bottomrule
    \end{tabular}
\end{table}

Sul primo frutto da raccogliere (Rank~1), entrambi i metodi ottengono lo 
stesso risultato complessivo, pari all'$81.0\%$ ($17$ scelte concordi su $21$). 
Andando però a esaminare le singole immagini, si notano delle differenze:
\begin{itemize}
    \item In $19$ scene su $21$ ($90.5\%$), la baseline e il planner ellittico 
    scelgono lo stesso identico frutto. In $16$ di questi casi la decisione 
    coincide con quella umana, mentre nelle restanti $3$ scene entrambi i metodi 
    scelgono un pomodoro diverso rispetto all'operatore. Come illustrato nel caso 
    di Figura~\ref{fig:soggettivita_gt}, in queste scene i frutti appartengono 
    allo stesso gruppo: l'algoritmo assegna il Rank~1 al pomodoro in alto 
    ($G = 0.533$) poiché risulta più isolato e centrale nel fotogramma, 
    mentre l'operatore umano ha preferito iniziare la raccolta dal frutto grande 
    in basso ($H_1$). Questa divergenza riflette la naturale soggettività 
    dell'annotatore umano rispetto al calcolo deterministico del punteggio geometrico.
\end{itemize}

\begin{figure}[H]
    \centering
    \includegraphics[width=\textwidth]{figures/fig4_soggettivita_gt.jpg}
    \caption[Confronto tra preferenza umana e ottimizzazione geometrica]
    {Esempio di divergenza tra annotazione umana e punteggio geometrico: 
    (a)~Ground Truth umano ($H_1, H_2, H_3$), che avvia la raccolta dal frutto 
    grande inferiore; (b)~Baseline a bounding box ($C_1, C_2, C_3$) e (c)~Planner 
    ellittico ($E_1, E_2, E_3$), che assegnano entrambi il Rank~1 al pomodoro 
    superiore isolato in virtù del maggiore punteggio geometrico ($G = 0.533$ 
    contro $0.490$).}
    \label{fig:soggettivita_gt}
\end{figure}

\begin{itemize}
    \item Nelle altre $2$ scene i due metodi prendono strade diverse. In 
    particolare, nel caso mostrato in Figura~\ref{fig:confronto_trittico}, la 
    bounding box seleziona per errore un frutto centrale coperto ($C_1$), 
    mentre l'ellisse individua correttamente il frutto libero sulla destra ($E_1$), 
    allineandosi con l'operatore umano ($H_1$). Nella seconda scena, al contrario, 
    è la baseline a coincidere con l'umano, mentre il nostro metodo sceglie un 
    frutto adiacente con un'area visibile leggermente maggiore.
\end{itemize}

Il vero vantaggio delle ellissi si osserva sulle scelte successive (Rank~2 e 
Rank~3). Man mano che la raccolta procede, la rimozione virtuale del primo 
pomodoro libera spazio e attiva il bonus di sblocco $\beta_{\text{unlock}}$ per i 
frutti retrostanti. Con i riquadri rigidi (BBox), gli angoli delle scatole 
continuano a sovrapporsi anche quando i frutti non si toccano, introducendo 
false occlusioni; di conseguenza, la concordanza con l'umano cala rapidamente 
al $38.5\%$ sul Rank~2 e al $20.0\%$ sul Rank~3. Con le ellissi, invece, la forma 
curva del pomodoro viene approssimata meglio: la concordanza resta al $46.2\%$ 
sul secondo frutto e raggiunge il $40.0\%$ sul terzo, raddoppiando il valore della 
baseline. Questo conferma che l'uso delle ellissi e l'aggiornamento dinamico delle 
occlusioni servono soprattutto a pianificare bene l'intera sequenza di prelievo, 
evitando errori a catena dopo la prima presa.


\subsection{Origine Geometrica delle Divergenze}
\label{subsec:origine_divergenze}

Come illustrato nel confronto qualitativo di Figura~\ref{fig:confronto_trittico}, 
le divergenze emerse tra i due approcci derivano da tre fattori 
intrinseci alla rappresentazione geometrica:

\begin{figure}[H]
    \centering
    \includegraphics[width=\textwidth]{figures/fig4_1_confronto_trittico.jpg}
    \caption[Confronto qualitativo delle strategie di picking su grappolo denso]{Confronto 
    qualitativo delle strategie di picking su un grappolo denso: 
    (a)~Ground Truth umano ($H_1, H_2, H_3$), con priorità sul frutto frontale destro; 
    (b)~Baseline a bounding box ($C_1, C_2, C_3$), che tenta erroneamente il prelievo 
    del frutto centrale occluso; 
    (c)~Planner ellittico proposto ($E_1, E_2, E_3$), che modella correttamente l'ingombro 
    e assegna il Rank~1 al pomodoro libero, in piena concordanza con l'operatore.}
    \label{fig:confronto_trittico}
\end{figure}

\begin{enumerate}
    \item \textbf{Eliminazione delle false occlusioni d'angolo}: Un riquadro 
    rettangolare allineato agli assi (AABB) racchiude il pomodoro includendo anche 
    i quattro angoli vuoti, che occupano circa il $21.5\%$ dell'intera scatola 
    (pari a un'area vuota di oltre il $27\%$ rispetto all'ingombro effettivo del 
    frutto, valore che cresce ulteriormente se il pomodoro è allungato e inclinato). 
    Quando due pomodori si trovano vicini in diagonale, gli angoli delle rispettive 
    scatole si sovrappongono pur senza alcun contatto reale tra i frutti. Questo 
    porta la baseline a rilevare una falsa occlusione e ad applicare la penalità 
    $\gamma_{\text{occ}} = 0.70$ a un frutto che in realtà è libero.
    \item \textbf{Rappresentazione dell'eccentricità e dell'orientamento}: 
    I pomodori presentano frequentemente forme ovoidali. L'ellisse orientata 
    cattura fedelmente l'angolo $\theta_i$ e i semiassi del frutto, 
    stimando con precisione l'effettiva area di contatto $\rho_{ij}$ tra corpi 
    adiacenti.
    \item \textbf{Pianificazione dinamica Pick \& Update}: La baseline valuta 
    la raggiungibilità in un singolo passaggio statico. Al contrario, la 
    pipeline proposta simula la rimozione del frutto selezionato e aggiorna le 
    occlusioni residue: l'assegnazione del bonus di sblocco 
    $\beta_{\text{unlock}} = 1.20$ promuove i frutti retrostanti liberati 
    dall'azione precedente, rispecchiando l'ordine di rimozione progressivo 
    adottato naturalmente dall'operatore umano.
\end{enumerate}

\section{Analisi dei Casi Studio e Discussione}
\label{sec:casi_studio_discussione}

L'analisi dei casi studio consente di osservare il comportamento dell'algoritmo 
su situazioni reali complesse, verificando da vicino come viene gestito lo 
sblocco dei frutti occlusi, quali fattori determinano la sequenza di raccolta 
e quali sono i principali limiti pratici del sistema.


\subsection{Dinamica di Sblocco delle Occlusioni}
\label{subsec:caso_studio_sblocco}

Il primo caso studio analizza il comportamento della pipeline decisionale sul 
fotogramma ad alta densità illustrato in Figura~\ref{fig:caso_sblocco}. La 
scena presenta un grappolo composto da tre pomodori maturi ($T_1, T_2, T_3$) 
e due frutti acerbi sottostanti:
\begin{itemize}
    \item $T_1$, posizionato inferiormente a coordinate $(245, 262)$, con 
    punteggio geometrico iniziale $G_1 = 0.482$;
    \item $T_2$, situato a coordinate $(257, 207)$ immediatamente alle spalle 
    di $T_1$, con punteggio geometrico di base $G_2 = 0.466$;
    \item $T_3$, collocato lateralmente a coordinate $(314, 179)$, privo di 
    ostacoli frontali, con punteggio geometrico $G_3 = 0.475$.
\end{itemize}

\begin{figure}[H]
    \centering
    \includegraphics[width=\textwidth]{figures/fig4_2_caso_studio_sblocco.jpg}
    \caption[Dinamica di sblocco su grappolo denso]{Dinamica di sblocco su un 
    grappolo denso con occlusioni multiple. 
    (a)~Ground Truth umano ($H_1, H_2, H_3$); 
    (b)~Baseline a bounding box ($C_1, C_2, C_3$); 
    (c)~Planner ellittico ($E_1, E_2, E_3$), in cui la rimozione di $E_1$ 
    sblocca $E_2$ assegnandogli il bonus $\beta_{\text{unlock}} = 1.20$ e 
    consentendogli di scavalcare il frutto laterale $E_3$.}
    \label{fig:caso_sblocco}
\end{figure}

Alla prima iterazione ($t = 1$), l'algoritmo valuta $T_1$ come frutto più 
accessibile e lo assegna al Rank~1 ($S_1 = 0.482$). Poiché $T_1$ si trova 
frontalmente rispetto a $T_2$, la simulazione del suo prelievo azzera la 
sovrapposizione reciproca ($\rho_{12} = 0$), liberando il target retrostante. 
Al passo successivo ($t = 2$), al pomodoro $T_2$ viene quindi applicato il 
bonus di sblocco $\beta_{\text{unlock}} = 1.20$, aggiornandone il punteggio a:
\begin{equation}
    S_2 = G_2 \times \beta_{\text{unlock}} = 0.466 \times 1.20 = 0.559
\end{equation}
Senza questo incentivo, $T_2$ ($0.466$) sarebbe rimasto subordinato a $T_3$ 
($0.475$), spingendo il pianificatore a selezionare il frutto laterale. Grazie 
al moltiplicatore dinamico, $T_2$ scavalca $T_3$ ($0.559 > 0.475$) e viene 
raccolto come Rank~2, lasciando $T_3$ al Rank~3.

Dal punto di vista pratico, questo comportamento è vantaggioso per due motivi:
\begin{itemize}
    \item \textbf{Sicurezza meccanica}: poiché $T_1$ si trova davanti a $T_2$, 
    tentare di raccogliere direttamente $T_2$ comporterebbe il rischio di urtare 
    e danneggiare il frutto antistante con le dita della pinza. Raccogliere 
    prima $T_1$ libera il corridoio di accesso frontale per il braccio.
    \item \textbf{Continuità sul grappolo locale}: all'inizio il pomodoro $T_2$ 
    aveva un punteggio più basso ($G_2 = 0.466$) rispetto al frutto laterale $T_3$ 
    ($G_3 = 0.475$) solo perché l'occlusione di $T_1$ ne nascondeva una parte, 
    riducendone l'area visibile alla telecamera. Senza il bonus di sblocco, dopo aver 
    tolto $T_1$ il robot sarebbe "saltato" sul frutto laterale $T_3$, lasciando 
    indietro $T_2$. Il moltiplicatore $\beta_{\text{unlock}} = 1.20$ compensa 
    questo svantaggio visivo iniziale, consentendo al manipolatore di completare 
    la raccolta dei pomodori contigui sullo stesso grappolo prima di spostarsi su 
    altri tralci.
\end{itemize}


\subsection{Inversione delle Priorità e Frutti Allungati}
\label{subsec:caso_studio_inversione}

Il secondo caso studio, illustrato in Figura~\ref{fig:caso_inversione}, 
mostra come l'uso delle ellissi risolva i limiti dei modelli rettangolari 
quando i pomodori hanno una forma allungata.

\begin{figure}[H]
    \centering
    \includegraphics[width=\textwidth]{figures/fig4_3_caso_studio_inversione.jpg}
    \caption[Inversione delle priorità di picking su grappolo verticale]{Inversione 
    delle priorità di picking su un grappolo verticale. (a)~Scelta umana di 
    riferimento ($H_1, H_2, H_3$). (b)~Baseline BBox ($C_1, C_2, C_3$), che 
    sceglie per errore il frutto centrale ($C_2$) prima di quello in basso 
    ($C_3$). (c)~Planner ellittico ($E_1, E_2, E_3$), che si adatta alla forma 
    allungata del pomodoro inferiore assegnandogli il Rank~2 ($E_2$), 
    rispecchiando esattamente la scelta dell'operatore umano.}
    \label{fig:caso_inversione}
\end{figure}

La scena mostra un grappolo verticale con tre pomodori in primo piano lungo il 
tralcio. Entrambi i metodi concordano nell'assegnare il Rank~1 al pomodoro in 
alto a coordinate $(147, 163)$, che è maturo e del tutto libero da ostacoli. 
La differenza fondamentale riguarda la scelta del secondo frutto da 
raccogliere, tra quello centrale a $(145, 254)$ e quello in basso a $(126, 331)$:
\begin{itemize}
    \item La baseline a Bounding Box seleziona il pomodoro centrale $(145, 254)$ 
    come Rank~2 ($C_2$) e lascia per ultimo quello in basso $(126, 331)$ al 
    Rank~3 ($C_3$);
    \item Il nostro planner a ellissi inverte l'ordine: raccoglie prima il 
    pomodoro in basso $(126, 331)$ come Rank~2 ($E_2$) e poi quello centrale 
    $(145, 254)$ come Rank~3 ($E_3$), esattamente nello stesso ordine scelto 
    dall'operatore umano ($H_1 \to H_2 \to H_3$).
\end{itemize}

Questo cambiamento è dovuto alla forma allungata del pomodoro inferiore, che 
ha un'altezza quasi doppia rispetto alla larghezza (semiassi di $30.7\text{ px}$ 
e $57.1\text{ px}$, con un rapporto di $1.86:1$). La baseline usa un riquadro 
rettangolare che allarga troppo l'area occupata ai lati, facendo credere al 
sistema che il pomodoro centrale sia sopra a quello inferiore. L'ellisse 
inclinata ($\theta \approx 3.2^\circ$), invece, segue con precisione la forma 
reale del frutto: l'algoritmo capisce quindi che il pomodoro in basso è libero 
e può essere raccolto subito, senza rischiare collisioni.

Oltre alla corretta stima delle sovrapposizioni, questo caso evidenzia un netto 
miglioramento nel calcolo del punteggio geometrico $G$. Nella baseline di 
riferimento, il perimetro del frutto viene calcolato direttamente sul contorno 
grezzo della maschera di segmentazione, che presenta bordi frastagliati e a 
gradino; questo fa aumentare artificialmente il perimetro misurato 
($P = 313.2\text{ px}$) e abbassa la circolarità a $C = 0.631$. Nel nostro 
approccio, l'uso dell'involucro convesso (\textit{Convex Hull}) prima del fitting 
dell'ellisse regolarizza i contorni, riducendo il perimetro a $P = 283.1\text{ px}$ 
e portando la circolarità a $C = 0.772$ ($+0.141$). Questo incremento 
migliora direttamente lo score geometrico $G$, fornendo una stima della 
raggiungibilità più realistica e stabile rispetto alla maschera grezza.


\subsection{Limiti Sistemici e Condizioni Critiche}
\label{subsec:limiti_sistemici}

L'analisi sperimentale ha evidenziato due limiti pratici della pipeline, 
che indicano le direzioni per i futuri sviluppi:
\begin{enumerate}
    \item \textbf{Zone di forte ombra e filtro di maturità}: La distinzione tra 
    pomodori maturi e acerbi si basa su soglie minime di colore e luminosità 
    nello spazio HSV ($S > 50$ e $V > 50$). Quando un frutto si trova sotto una 
    fitta ombra fogliare, molti dei suoi pixel risultano troppo bui per 
    superare la soglia. La forte riduzione dei pixel analizzabili rende il 
    calcolo della maturità più fragile ed esposto a disturbi dello sfondo. In 
    futuro, questo limite potrà essere superato adottando soglie che si adattano 
    automaticamente alla luce della scena, oppure addestrando una piccola rete 
    neurale dedicata unicamente a riconoscere la maturazione.
    \item \textbf{Frutti completamente nascosti e visione da singolo punto 
    di vista}: L'algoritmo elabora una sola inquadratura fotografica fissa. 
    Se un pomodoro maturo è interamente coperto da foglie o posizionato dietro 
    al fusto della pianta, la telecamera non può rilevarlo in alcun modo. 
    Questo evidenzia il limite della visione statica e la necessità di passare 
    a strategie di \textit{Active Vision}: muovendo la telecamera con il braccio 
    robotico attorno al grappolo, il sistema potrà esplorare la pianta da diverse 
    angolazioni e individuare anche i frutti inizialmente non visibili.
\end{enumerate}
